\begin{corollary}[Stationary class under norm pullback]\label{mod:parameters-stationaryclassundernorm}
Let $K/F$ be a finite extension of nonarchimedean local fields.  Let
$\chi_F$ and $\chi_K$ be quasi-character data with their exact, separately
stored conductors
\[
 m_F=m_F(\chi_F),\qquad m_K=m_K(\chi_K),
\]
and let $\psi_F$ and $\psi_K$ be exact additive-character data.  Assume the
manuscript's pullback normalizations
\[
 \chi_K=\chi_F\circ N_{K/F},\qquad
 \psi_K=\psi_F\circ\operatorname{Tr}_{K/F}.
\]
Choose natural numbers $d_F,d_K,\varepsilon_F,\varepsilon_K$ and a
\texttt{NormPolynomialPrecision} certificate $h$ for
\[
 m_F=2d_F+\varepsilon_F,qquad
 m_K=2d_K+\varepsilon_K,qquad
 \varepsilon_F,\varepsilon_K\leq1.
\]
The certificate uses the actual conductor $m_K$ of $\chi_K$, not $m_F$, as
its upstairs conductor.  It also proves $m_F,m_K>1$ and records separately
\[
 N(U_K^{d_K+\varepsilon_K})\subseteq U_F^{d_F+\varepsilon_F},qquad
 N(U_K^{m_K})\subseteq U_F^{m_F},qquad
 N(U_K^{d_K})\subseteq U_F^{d_F}.
\]
The first two inclusions induce the additive truncated norm
\[
 P_h:\mathfrak p_K^{d_K+\varepsilon_K}/\mathfrak p_K^{m_K}
 \longrightarrow
 \mathfrak p_F^{d_F+\varepsilon_F}/\mathfrak p_F^{m_F},
 \qquad [y]\longmapsto[N_{K/F}(1+y)-1].
\]

Choose the ordered denominator pair $(\gamma_F,\gamma_K)$ with
\[
 \operatorname{ord}_F(\gamma_F)=m_F+n_F(\psi_F),qquad
 \operatorname{ord}_K(\gamma_K)=m_K+n_K(\psi_K).
\]
Lean explicitly transports the minimal-depth Lamprecht stationary class to
the literal coefficient quotient and defines
\[
 \beta_F=\operatorname{St}_{\gamma_F,d_F+\varepsilon_F}(\chi_F)
   \in C_F:=\mathcal O_F/\mathfrak p_F^{d_F},
 \qquad
 \beta_K=\operatorname{St}_{\gamma_K,d_K+\varepsilon_K}(\chi_K)
   \in C_K:=\mathcal O_K/\mathfrak p_K^{d_K}.
\]
This transport is an additive equivalence which keeps every supplied
integral representative and uses
$m_E-(d_E+\varepsilon_E)=d_E$; it makes no representative choice.

For this same ordered denominator pair, the denominator-sensitive adjoint has
the reverse direction
\[
 P^*_{h;\gamma_F,\gamma_K}:C_F\longrightarrow C_K.
\]
Then Lean proves the quotient-class equality
\[
 \boxed{\ \beta_K=P^*_{h;\gamma_F,\gamma_K}(\beta_F)\ }
 \qquad\text{in }\mathcal O_K/\mathfrak p_K^{d_K}.
\]
Indeed, pullback along the truncated norm identifies the downstairs and
upstairs stationary linearization characters.  Equivalently, the adjoint
image satisfies, for every
$[y]\in\mathfrak p_K^{d_K+\varepsilon_K}/\mathfrak p_K^{m_K}$,
\[
 \chi_K(1+y)=
 \Phi_{K,\gamma_K}
   \bigl(P^*_{h;\gamma_F,\gamma_K}(\beta_F)\bigr)([y]),
\]
where $\Phi_{K,\gamma_K}$ is the quotient pairing from
Definition~\ref{mod:parameters-adjoint}.
The proof uses the exact norm identity represented by $P_h$, not a replacement
of norm by trace.  It then transports
\texttt{stationaryNumeratorClass\_unique} from
Lemma~\ref{mod:lamprecht-stationaryclass} at the actual upstairs conductor
$m_K$ and the exact depth $d_K+\varepsilon_K$.

All field elements in the representative statements are separately supplied.
If $c,c'\in\mathcal O_F$ both represent $\beta_F$, and $b,b'\in\mathcal O_K$
represent the corresponding adjoint images, then
\[
 b\equiv b'\pmod{\mathfrak p_K^{d_K}}.
\]
Every such $b$ satisfies the displayed upstairs linearization.  If
$y\equiv y'\pmod{\mathfrak p_K^{m_K}}$, both the downstairs norm-polynomial
phase and the upstairs phase are unchanged.  Thus none of these assertions
promotes a quotient class to a canonical field representative.

The class is compatible with every deeper admissible stationary depth: the
canonical coefficient-quotient restriction of the minimal-depth class is exactly the
stationary class constructed at that deeper depth.  It is also compatible
with admissible denominator changes.  If $\gamma_E'$ is another admissible
denominator, let
\[
 S_E^{\gamma_E,\gamma_E'}([c])=
   [\,(\gamma_E'/\gamma_E)c\,]
   \quad\text{on }\mathcal O_E/\mathfrak p_E^{d_E}.
\]
Then the stationary class for $\gamma_E'$ is
$S_E^{\gamma_E,\gamma_E'}$ of the class for $\gamma_E$, and Lean proves on
the downstairs stationary class the denominator-sensitive square
\[
 P^*_{h;\gamma_F',\gamma_K'}
   \bigl(S_F^{\gamma_F,\gamma_F'}(\beta_F)\bigr)
 =S_K^{\gamma_K,\gamma_K'}
   \bigl(P^*_{h;\gamma_F,\gamma_K}(\beta_F)\bigr).
\]

Finally, the same precision certificate proves $m_F,m_K>1$.  Hence this
corollary constructs no stationary class when either actual conductor is zero
or one; those endpoint cases remain outside its hypotheses.

\LeanFile{LanglandsFirstMainLemma/Parameters/StationaryClassUnderNorm.lean}
\lean{LanglandsFirstMainLemma.stationaryDepthOfConductorDecomposition}
\lean{LanglandsFirstMainLemma.stationaryCoefficientLamprechtEquivAtConductor}
\lean{LanglandsFirstMainLemma.stationaryCoefficientLamprechtEquivAtConductor_mk}
\lean{LanglandsFirstMainLemma.stationaryCoefficientClass}
\lean{LanglandsFirstMainLemma.stationaryCoefficientClass_toLamprecht}
\lean{LanglandsFirstMainLemma.stationaryPairingLeft_eq_lamprechtPairing}
\lean{LanglandsFirstMainLemma.IsStationaryCoefficientClass}
\lean{LanglandsFirstMainLemma.stationaryCoefficientClass_isStationary}
\lean{LanglandsFirstMainLemma.stationaryCoefficientClass_unique}
\lean{LanglandsFirstMainLemma.stationaryCoefficientClass_existsUnique}
\lean{LanglandsFirstMainLemma.latticeQuotientMk_eq_stationaryCoefficientClass_iff}
\lean{LanglandsFirstMainLemma.stationaryCoefficientClass_restrict}
\lean{LanglandsFirstMainLemma.stationaryLinearizationCharacter_compNorm}
\lean{LanglandsFirstMainLemma.normPolynomialAdjoint_stationaryClass_isStationary}
\lean{LanglandsFirstMainLemma.stationaryClass_compNorm}
\lean{LanglandsFirstMainLemma.stationaryClass_compNorm_representative_linearization}
\lean{LanglandsFirstMainLemma.stationaryClass_compNorm_representatives_independent}
\lean{LanglandsFirstMainLemma.stationaryClass_compNorm_variable_representative_independent}
\lean{LanglandsFirstMainLemma.stationaryCoefficientClass_changeDenominator}
\lean{LanglandsFirstMainLemma.stationaryClass_compNorm_changeDenominator}
\lean{LanglandsFirstMainLemma.stationaryClass_compNorm_conductors_gt_one}
\leanok
\SourceText{Corollary~\ref{cor:stationary-class-under-norm}.}
\uses{mod:parameters-adjoint,mod:lamprecht-stationaryclass,mod:lamprecht-formula}
\FormalizationContract The Lean theorem is parameterized by the three
certified norm-filtration inclusions rather than by a ramification case split.
Its coefficient quotients have exactly the depths $d_F$ and $d_K$ from the
separate exact conductor decompositions, and its variable quotients have
exactly the depths $d_F+\varepsilon_F$ and $d_K+\varepsilon_K$.  The truncated
norm points from upstairs variables to downstairs variables, while the
adjoint points from downstairs coefficient classes to upstairs coefficient
classes and retains both ordered denominators.  Stationary uniqueness is used
only after proving the upstairs linearization at the actual upstairs depth.
All principal equalities are finite lattice-quotient equalities; all
representatives are explicit hypotheses, and no stationary object is defined
at conductor zero or one.
\end{corollary}
